\section{Method}
\label{sec:method}

This section summarises the ONE-NAS algorithm \cite{lyu2023onenas},
defines the two prediction rules compared in this study, and describes
the deployment architecture in which search and inference are executed
on separate hardware. Configuration values are deferred to
Section~\ref{sec:setup}.

\subsection{Online Neuroevolution}

ONE-NAS maintains a population of RNNs and evolves both their
architectures and their weights as data arrive. Each individual, termed
a genome, is a directed graph in which hidden nodes are instantiated
from a library of recurrent cell types. Evolution begins from a minimal
genome containing only input-to-output connections. Offspring are
produced by mutation (insertion or deletion of nodes and edges,
insertion of recurrent edges, and change of cell type) and by crossover
between two parents. Weights are inherited from the parent for every
retained component and randomly initialised only for newly introduced
components. This Lamarckian inheritance reduces the training required
per offspring to a small number of backpropagation epochs and is a
prerequisite for operating within the data cadence.

The population is partitioned into $N$ islands. Each island maintains an
elite set of fixed size, generates a fixed number of offspring per
generation from that set, and admits genetic material from other
islands only through inter-island crossover at a low rate. Islands
consequently converge to distinct regions of the architecture space. To
counteract stagnation, at a fixed interval the island whose best elite
has the worst fitness is cleared and re-initialised from the global
best genome.

\subsection{Online Protocol}

The input stream is partitioned into windows of fixed length in trading
days, and one generation is executed per window. Let $\mathcal{H}_t$
denote the history available at generation $t$. In generation $t$, each
island generates offspring from its elites; offspring are trained on
worker processes using subsequences sampled from $\mathcal{H}_t$; all
genomes, elite and offspring, are evaluated by mean squared error on
the most recent windows of $\mathcal{H}_t$, which are excluded from
training; and each island retains the top-ranked genomes as its elite
set. The window revealed during generation $t$ is then appended to form
$\mathcal{H}_{t+1}$. Because fitness is measured on recent data while
training draws on the full history, the population tracks the current
regime without discarding earlier ones. Genome generation is performed
by a single master process and training is distributed asynchronously
across workers, so that the wall-clock cost of a generation is governed
by the number of workers rather than the population size.

\subsection{Prediction Rules}
\label{sec:prediction_rules}

At the end of generation $t$, the algorithm exposes the elite set of
each island and the genome of best fitness across all islands. The
selection of genomes used to predict window $t{+}1$ is a scoring-time
decision that does not affect the search. Two rules are compared.

Under the \emph{single champion} rule, the global best genome alone
produces the prediction. This is the rule used in all prior ONE-NAS
work.

Under the \emph{island-champion ensemble} rule, the best elite of each
island produces a prediction, giving $N$ predicted cross-sections. Each
is converted to ranks, and the ranks are averaged across members. Rank
aggregation is used in place of output averaging because the members
are independently trained networks with non-comparable output scales,
and because every trading rule considered here depends only on the
ordering of the prediction.

Both rules are evaluated on the same runs. Any difference between them
is therefore attributable to the use of the evolved population rather
than to its construction.

\subsection{Deployment Architecture}
\label{sec:deployment}

Search and inference are executed on separate machines, reflecting the
asymmetry in their computational requirements: population training is
parallel and expensive, whereas evaluating the resulting networks is
not. Figure~\ref{fig:deployment} depicts the pipeline.

\paragraph{Search host.} ONE-NAS is executed on a cloud instance or a
workstation with sufficient cores to train a generation's offspring in
parallel, using one master process and a pool of training workers. The
host stores the price history, advances the window clock, and performs
all evolutionary operations.

\paragraph{Endpoint.} Inference is performed on a Raspberry Pi,
selected as a deliberately constrained target with substantially less
compute and memory than a typical consumer laptop or office desktop.
Demonstrating that inference meets the data cadence on this device
establishes a lower bound on the hardware required at the point of use.
The endpoint stores only the most recently received genomes and the
current day's feature vectors; no training is performed on the device.

\paragraph{Model hand-off.} Transfer of genomes from host to endpoint is
a step of every generation. Upon completion of selection in generation
$t$, the master serialises the genomes designated to predict window
$t{+}1$ (the global best genome, or the $N$ island champions, according
to the prediction rule) into a binary encoding of graph structure, cell
types, and weights, and transmits the file to the endpoint over TCP/IP.
The endpoint deserialises the networks and produces a prediction for
each day of window $t{+}1$ as its features become available, applying
the rank aggregation of Section~\ref{sec:prediction_rules} when an
ensemble is received. Concurrently, the host generates and trains
generation $t{+}1$. Prediction and search thus overlap in time, as in
the original ONE-NAS design, but are decoupled across hardware.

Two properties of this architecture are relevant to deployment. First,
inference cost scales only with $N$ and the number of securities, so
increasing the island count does not alter the endpoint's hardware
requirement. Second, the endpoint is stateless with respect to the
search: in the event of host unavailability, it continues to predict
with the most recently received genomes and resumes upon the next
hand-off.